\section{Convergence: Law of Large Numbers}
\label{sec: law of large numbers}

The difference between a ``weak law'' and ``strong law'' is the convergence type.
A ``weak law'' only requires convergence in probability, while a ``strong law''
requires almost sure convergence. The typical strategy to obtain strong laws is
to obtain quantitative convergence results in expectation (e.g. in \(L^p\)), translate them
to quantitative convergence result in probability using the Markov inequality
and deduce almost sure convergence using the Borel-Cantelli lemmas from
the rate of convergence, that is
\[\begin{aligned}
    \substack{\text{Qunatitative}\\ X_n \overset{L^p}\to X } \quad
    \overset{\text{Markov ineq.}}\implies \quad \substack{\text{Quantitative}\\ X_n \overset{p}\to X}\quad
    \overset{\text{Borel-Cantelli}}\implies \quad X_n \overset{\as}\to X.
\end{aligned}
\]
The reason for this strategy is that it is often easiest to calculate
expectations.  For example \(\norm{X - \E{X}}_{L^2} = \sqrt{\var{X}}\). 
In this section we illustrate this strategy by implementing it for the
law of large numbers. It turns out that the natural domain to work in
is a Hilbert space \((\hilbert, \scp{\cdot, \cdot})\) as it allows
the definition of the covariance and variance in a natural way.

\begin{remark}[Expectation in multiple dimensions]
    \label{rem: expectation in multiple dimensions}
    In the discrete case, the expectation is
    \[
        \E{X} = \sum_{x} x \prob{X = x}
    \]
    Observe the key for this to work is that the \(x\) are summable and scalable
    by a scalar probability. This only requires a vector space structure.  
    In general the expectation is defined via integrals -- specifically the
    ``Bochner integral'' in infinite dimension. We will not go into the measure
    theoretical details here. Simply understand that the expectation is a linear
    and therefore an ``entrywise'' operation acting on the basis due to
    \[
        \E{X}
        = \E[\Big]{\begin{pmatrix}
            X^{(1)} \\
            \vdots \\
            X^{(\dims)}
        \end{pmatrix}}
        = \E[\Big]{\sum_{i=1}^\dims X^{(i)} e_i}
        \overset{\text{linear}}= \sum_{i=1}^\dims \E{X^{(i)}} e_i
        = \begin{pmatrix}
            \E{X^{(1)}} \\
            \vdots \\
            \E{X^{(\dims)}}
        \end{pmatrix}.
    \]
\end{remark}

\begin{definition}[Covariance]
    \label{def: covariance and variance definition}
    Let \(X, Y \in \cL^2(\hilbert)\), then their \textbf{covariance}
    and \textbf{variance} are defined as
    \[
        \cov{X}{Y} \coloneq \E[\Big]{\scp[\big]{X-\E{X}, Y-\E{Y}}}
        \quad \text{and} \quad
        \var{X} \coloneq \cov{X}{X}.
    \]
\end{definition}

\begin{remark}
Note that \(\covSym\) is an inner product and we will later define
a covariance operator \(\CovSym\) that is an outer product such that \(\covSym\)
becomes the trace of \(\CovSym\).
\end{remark}
The variance of sums or random variables is straightforward to calculate:

\begin{mdframed}[innertopmargin=0pt]
\begin{lemma}[Bienaymé]
    \label{lem: Bienaymé}
    Let \(X_1, \dots, X_n\) be random variables in \(\cL^2(\hilbert)\). Then
    \[
        \var[\Big]{\sum_{i=1}^n X_i} = \sum_{i=1}^n \var{X_i} + \sum_{\substack{i,j=1\\i \neq j}}^n \cov{X_i}{X_j}.
    \]
    If the random variables are uncorrelated (i.e.\ \(\cov{X_i}{X_j} = 0\) for \(i\neq j\)), then
    % for the mean \(\mean{X}_n \coloneq \frac1n \sum_{i=1}^n X_i\) we have
    \[
        % \var[\Big]{\sum_{i=1}^n X_i} = \sum_{i=1}^n \var{X_i}
        % \quad \text{and} \quad
        \var[\Big]{\frac1n \sum_{i=1}^n X_i} = \frac1{n^2} \sum_{i=1}^n \var{X_i}.
    \]
\end{lemma}
\end{mdframed}

\begin{proof}
    Let \(Y_i \coloneq X_i - \E{X_i}\). Then by definition of variance and covariance
    \[
        \var[\Big]{\sum_{i=1}^n X_i}
        % = \E[\Big]{\norm[\Big]{\sum_{i=1}^n Y_i}^2}
        = \sum_{i,j=1}^n \E[\big]{\scp{Y_i, Y_j}}
        = \sum_{i,j=1}^n \cov{X_i}{X_j}.
    \]
    With \(\var{X_i} = \cov{X_i}{X_i}\) the first result follows from a split of the sum.
    The second result follows directly from the first and
    \(
        \var{aX} = a^2 \var{X}
    \) for any \(a\in \real\).
\end{proof}

If we do not care about the speed of convergence, then an immediate corollary of
Bienaymé is the following weak law of large numbers.

\begin{corollary}[Khinchin's weak law of large numbers]
    \label{cor: weak law of large numbers}
    Let \((X_n)_{n\in \nat}\) be a sequence of uncorrelated random variables
    with identical mean \(\mu= \E{X_n}\) and \(\frac1{n^2} \sum_{i=1}^n \var{X_i} \to 0\).
    Then
    \[
        \mean{X}_n\coloneq \frac1n \sum_{i=1}^n X_i \overset{p}\to \mu.
    \]
\end{corollary}
\begin{proof}
    By the Chebyshev inequality (Corollary \ref{cor: markov inequality Lp})
    we get
    \begin{equation*}
        \begin{aligned}
        \prob[\big]{\norm{\mean{X}_n - \mu} > \epsilon}
        &\le \tfrac1{\epsilon^2}\|\mean{X}_n - \mu\|_{L^2}^2
        \\
        &= \frac1{\epsilon^2} \var[\big]{\mean{X}_n}
        \overset{\text{Bienaymé}}= \frac1{\epsilon^2}\Bigl(\frac1{n^2} \sum_{i=1}^n \var{X_i}\Bigr).
        \end{aligned}
    \end{equation*}
    The assumption \(\frac1{n^2} \sum_{i=1}^n \var{X_i} \to 0\) therefore
    yields the claim.
\end{proof}

From the proof it is clear that the speed of convergence is determined by the
speed of convergence of \(\frac1{n^2} \sum_{i=1}^n \var{X_i}\).
If the variances are bounded, that is
\(\sup_{i}\var{X_i} \le \sigma^2 < \infty\),
then
\begin{equation}
    \label{eq: rate of convergence for bounded variance}    
    \prob[\big]{\norm{\mean{X}_n - \mu} > \epsilon}
    \le \frac1{\epsilon^2}\frac{\sigma^2}{n} \in \bigO(\tfrac1n).
\end{equation}
If the variances \(\var{X_i}\) are increasing in \(i\), then one
may still obtain (slower) convergence. But even the convergence speed for the
bounded case is too slow to deduce almost sure convergence directly as 
as the harmonic series diverges (cf.\ Corollary \ref{cor: characterization of
almost sure convergence}). The problem is essentially that Markov's inequality
is too lose. To prove almost sure convergence -- that is a strong
law of large numbers, a different approach is needed. There are
two common approaches:
\begin{enumerate}
    \item (Shortcut): \textbf{Bound the \(L^4\) norm} and use the Chebyshev inequality
    for \(p=4\). This results in a convergence speed on the order of
    \(\bigO(\frac1{n^2})\) which is summable and therefore sufficient for almost
    sure convergence by Corollary \ref{cor: characterization of almost sure
    convergence}.  However to bound the \(L^4\) norm it is necessary to assume
    the variables \(X_n\) have finite fourth moments \(\E{X_n^4} < \infty\) and
    assume independence such that \(\E{X_i X_j X_k X_l}\) factorizes when the
    indices are different (see Exercise \ref{ex: strong law of large numbers in
    L4}).

    \item The most general forms of the strong law of large numbers, is
    obtained by the observation that we get almost sure convergence
    on a \textbf{subsequence} (cf.\ Lemma \ref{lem: subsequences converge almost
    surely}). The second step is to control the elements between the
    converging subsequence.
\end{enumerate}

\begin{mdframed}[innertopmargin=0pt]%
\begin{theorem}[Strong law of large numbers]
    \label{thm: strong law of large numbers}
    Let \(\hilbert\) be a Hilbertspace and \((X_n)_{n\in \nat}\) a sequence
    of \underline{uncorrelated} random variables in \(\cL^2(\hilbert)\), with identical mean
    \(\mu = \E{X_n}\) and bounded variances \(\var{X_n} \le \sigma^2 < \infty\).
    Then
    \[
        \mean{X}_n \coloneq \frac1n \sum_{i=1}^n X_i \overset{\as}\to \mu.
    \]
\end{theorem}
\end{mdframed}

\begin{proof}[Proof (by Michael J Neely \citep{neelyInfinitelyOftenProbability2021})]
    Since the variance is bounded, we have seen in \eqref{eq: rate of convergence for bounded variance}
    that  \(\prob[\big]{\norm{\mean{X}_n - \mu} > \epsilon} \le
    \frac{\sigma^2}{n\epsilon^2}\). This implies \(\mean{X}_{n^2}
    \overset{\as}\to \mu\) by Corollary \ref{cor: characterization of almost sure
    convergence}
    since 
    \[
        \sum_{n=1}^\infty \prob[\big]{\norm{\mean{X}_{n^2}-\mu} \ge \epsilon}
        \le \frac{\sigma^2}{\epsilon^2}\sum_{n=1}^\infty \frac1{n^2} < \infty.
    \]
    That is we choose the  subsequence indexed by \(n^2\). What is left to prove
    is that \(X_n \overset{\as}\to 0\) for any \(n\in \nat\) that is not a
    perfect square. 
    Observe that by defining \(\tilde{X}_n \coloneq X_n - \mu\) we can
    assume without loss of generality that \(\mu=0\), which we will do from now on.
    With \(k_n \coloneq \lfloor \sqrt{n} \rfloor\) the closest 
    squares surrounding \(n\) are
    \begin{equation}
        \label{eq: k definition}    
        k_n^2 \le n < (k_n+1)^2.
    \end{equation}
    then we have
    \[
        \mean{X}_n = \underbrace{\frac1n \sum_{i=1}^{k_n^2} X_i}_{S_n}
        + \underbrace{\frac1n \sum_{i=k_n^2+1}^{n} X_i}_{\Delta_n}.
    \]
    We will show that \(S_n \overset{\as}\to 0\) and \(\Delta_n \overset{\as}\to
    0\) separately and conclude by continuous mapping that \(\mean{X}_n \overset{\as}\to
    0\). For the first term observe that by \eqref{eq: k definition} we have
    \[
        1 \ge
        \frac{k_n^2}{n}
        \ge \frac{k_n^2}{(k_n+1)^2} = \Bigl(\frac{1}{1 + \frac{1}{k_n}}\Bigr)^2
        \to 1.
    \]
    This implies \(\frac{k_n^2}{n} \to 1\). Since \(\mean{X}_{n^2} \overset{\as}\to 0\)
    implies \(\mean{X}_{k_n^2} \overset{\as}\to 0\)
    we have by continuous mapping
    \[
        S_n = \frac{k_n^2}{n} \mean{X}_{k_n^2} \overset{\as}\to 0.
    \]
    To show that the remaining term \(\Delta_n\) converges almost surely to zero we
    use the Chebyshev inequality with \(p=2\) again. To this end
    we bound the \(L^2\) norm of \(\Delta_n\). Observe that since the \(X_i\) are
    uncorrelated we have
    \[\begin{aligned}
        \var{\Delta_n}
        = \frac1{n^2} \sum_{i=k_n^2+1}^{n} \var{X_i}
        &\le \sigma^2 \frac{n - k_n^2}{n^2}
        \\
        \overset{\eqref{eq: k definition}}&\le \sigma^2 \frac{(k_n+1)^2 - k_n^2}{k_n^4}
        = \sigma^2 \frac{2k_n + 1}{k_n^4}
        \\
        \overset{k_n \ge 1}&\le \frac{3\sigma^2}{k_n^3}.
    \end{aligned}
    \]
    Consequently, by the Chebyshev inequality we have for every \(\epsilon > 0\)
    \[
        \sum_{n=1}^\infty \prob[\big]{|\Delta_n| \ge \epsilon}
        \le \sum_{n=1}^\infty \frac{\var{\Delta_n}}{\epsilon^2}
        \le \frac{3\sigma^2}{\epsilon^2}\sum_{n=1}^\infty \frac1{k_n^3}.
    \]
    Now we are finished by the characterization of almost sure convergence
    (Corollary \ref{cor: characterization of almost sure convergence}) if can
    show that this series converges. To this end observe that by definition
    of \(k_n\) we have \(k_n > \sqrt{n} -1\) and therefore
    \begin{align*}
        \sum_{n=3}^\infty \frac1{k_n^3}
        &\le \sum_{n=3}^\infty \frac1{(\sqrt{n} -1)^3}
        \\
        &\le \sum_{n=3}^\infty \int_{n-1}^n \frac1{(\sqrt{x} -1)^3} dx
        \\
        &= \int_2^\infty \frac1{(\sqrt{x} -1)^3} dx
        \\
        \overset{y=\sqrt{x}-1}&= \int_{\sqrt{2}-1}^\infty \frac{2(y+1)}{y^3} dy
        \le C \int_{\sqrt{2}-1}^\infty \frac1{y^2} dy 
        < \infty.
        \qedhere
    \end{align*}
\end{proof}

In Theorem \ref{thm: strong law of large numbers} we still assumed that the
random variables are in \(L^2\). And if you want to work with
``uncorrelated'' random variables this cannot be avoided as the covariance
is otherwise undefined. However, if you are willing to assume pairwise
independence, then the strong law of large numbers also holds under the weaker
assumption that the random variables are in \(L^1\).

\begin{theorem}[Etemadi's Strong law of large numbers]
    Let \((X_n)_{n\in \nat}\) be a sequence of identically distributed,
    \underline{pairwise} independent random variables
    in \(\cL^1\). Then \(\mean{X}_n \overset{\as}\to \E{X_1}\).
\end{theorem}
\begin{proof}[Proofsketch]
    Define the truncated random variables \(Y_n \coloneq \ind{|X_n| \le n}
    X_n\). Then the variance of every \(Y_n\) is
    finite although the variance is increasing in \(n\).
    Using Borel-Cantelli it is possible to show that eventually all \(Y_n\) are
    equal to \(X_n\). Then one proves the strong law of large numbers for the
    \(Y_n\). For this it is used that functions of pairwise independent random
    variables are also pairwise independent and consequently the \(Y_n\) are
    uncorrelated. We can unfortunately not
    apply Theorem \ref{thm: strong law of
    large numbers} since it requires a uniform bound on the variances and
    the variances are increasing here. For details see e.g.\
    \citep[Thm.~5.17]{klenkeProbabilityTheoryComprehensive2014}.
\end{proof}

\subsection*{Exercises}

\begin{exercise}[Uncorrelated vs independent \Coffeecup\Coffeecup]
   Let \(X\) and \(Y\) be two random variables in \(\real\). Prove that
   the following statements are equivalent:
   \begin{enumerate}[label=(\roman*),noitemsep]
        \item\label{it: independent} \(X\) and \(Y\) are independent, that is
        \[
            \prob{X \in A, Y \in B} = \prob{X \in A}\prob{Y \in B} \qquad \forall A, B
        \]
        \item\label{it: measurable functions} For all measurable functions \(f\) and \(g\) such that \(f(X)\)
        and \(g(X)\) are integrable we have
        \[
            \E{f(X)g(Y)} = \E{f(X)}\E{g(Y)}.
        \]
        \item\label{it: bounded measurable functions} For all \underline{bounded}, measurable functions \(f\) and \(g\)
        \[
            \E{f(X)g(Y)} = \E{f(X)}\E{g(Y)}.
        \]
   \end{enumerate}
   \textbf{Hint: } Use Theorem 102 from Ivan's lecture notes.
   \begin{solution}
        We have \ref{it: independent} \(\Rightarrow\) \ref{it: measurable functions} by Theorem 102 from Ivan's lecture.
        The implication \ref{it: measurable functions} \(\Rightarrow\) \ref{it: bounded measurable functions} is trivial.
        For the implication \ref{it: bounded measurable functions} \(\Rightarrow\) \ref{it: independent} observe that for any \(A, B\) we have
        that \(\ind{A}\) and \(\ind{B}\) are bounded measurable functions.
   \end{solution}
    
   \noindent Further prove the following statements
    \begin{enumerate}[label=(\alph*),noitemsep]
        \item If \(X\) and \(Y\) are independent, they are uncorrelated, that is
        \[
            \cov{X}{Y} = 0.
        \]
        \begin{solution}
            We have that
            \begin{align*}
                \cov{X}{Y}
                &= \E{(X-\E{X})(Y-\E{Y})}
                \\
                &= \E{XY} - \E{X}\E{Y}
                \overset{\text{indep.}}=
                \E{X}\E{Y} - \E{X}\E{Y} = 0.
                \qedhere
            \end{align*}
        \end{solution}

        \item Let \(\prob{X=-1} = \prob{X=0} = \prob{X=1} = \frac13\) and \(Y =
        X^2\). Prove that \(X\) and \(Y\) are uncorrelated but not independent.
        \begin{solution}
            Since \(\E{X} = 0\) we have that
            \begin{align*}
                \cov{X}{Y}
                &= \E{XY} - \E{X}\E{Y}
                \\
                &= \E{X^3} - \E{X}\E{X^2}
                = \E{X} - \E{X}\E{X^2}
                = 0.
            \end{align*}
            However, 
            \[
                \prob{X=0, Y=1} =0 \neq \tfrac13 \tfrac23 = \prob{X=0}\prob{Y=1}.
                \qedhere
            \]
        \end{solution}

        \item If \(X\) and \(Y\) are independent, then \(h_1(X)\) and \(h_2(Y)\) are independent
        for all measurable functions \(h_1\colon \real \to \real\) and \(h_2\colon \real \to \real\).
        \begin{solution}
            Since \(f\circ h_1\) and \(g \circ h_2\) are measurable for all
            measurable \(f\colon \real \to \real\) and \(g\colon \real \to \real\), we have by \ref{it: measurable
            functions} that
            \begin{align*}
                \E{f(h_1(X))g(h_2(Y))}
                &= \E{(f\circ h_1)(X)(g\circ h_2)(Y)}
                \\
                &= \E{f(h_1(X))}\E{g(h_2(Y))}.
            \end{align*}
            This proves independence by \ref{it: measurable functions}. 
        \end{solution}

        \item\label{it: all transformations uncorrelated implies independence}
        If \(f(X)\) and \(g(Y)\) are uncorrelated for all bounded
        measurable functions \(f\) and \(g\), then \(X\) and \(Y\) are
        independent.
        \begin{solution}
            We have
            \[
                0 = \cov{f(X)}{g(Y)} = \E{f(X),g(Y)} - \E{f(X)}\E{g(Y)}
            \]
            and consequently \(\E{f(X)g(Y)} = \E{f(X)}\E{g(Y)}\) for all
            bounded measurable functions \(f\) and \(g\).  This implies
            independence by \ref{it: bounded measurable functions}.
        \end{solution}
    \end{enumerate}

    \begin{remark}[Interpretation]
        What this essentially means is the following:
        \begin{enumerate}
            \item Uncorrelated is weaker than independence.
            \item \(f(X)\) and \(f(Y)\) are independent if \(X\) and \(Y\)
            are independent, that is independence is \emph{stable under transformation}.

            \item In contrast, \(X\) and \(Y\) being uncorrelated is \emph{not} stable under transformation.
            In fact, if it were true that \(f(X)\) and \(g(Y)\) is uncorrelated
            for all bounded measurable functions \(f\) and \(g\), then \(X\) and
            \(Y\) would be independent by \ref{it: all transformations
            uncorrelated implies independence}. Independence therefore gains the
            interpretation of ``all transformations are uncorrelated''.
        \end{enumerate}

    \end{remark}
\end{exercise}

\begin{exercise}[Mean of indicators \Coffeecup]
    \label{ex: mean of indicators}
    Let \((X_n)_{n\in \nat}\) be a sequence of independent and identically
    distributed random variables in \(\real\) with distribution function \(F\).
    \begin{enumerate}[label=(\alph*),noitemsep]
        \item Let \(A\) be a measurable set and \(Y_i \coloneq \ind{X_i \in A}\).
        Prove that
        \[
            \frac1n \sum_{i=1}^n Y_i \overset{\as}\to \prob{X_1 \in A}
            \qquad (n \to \infty).
        \]
        \begin{solution}
            This is the strong law of large numbers for the independent (and
            thereby uncorrelated), identically distributed \(Y_i\), since
            \[
                \E{Y_i} = \E{\ind{X_i \in A}} = \prob{X_1 \in A}.
                \qedhere
            \]
        \end{solution}

        \item Define the empirical distribution function
        \[
            F_n(x) \coloneq \frac1n \sum_{i=1}^n \ind{X_i \le x}.
        \]
        Prove that for any \(x_1, \dots, x_k \in \real\)
        \[
            (F_n(x_1), \dots, F_n(x_k)) \overset{\as}\to (F(x_1), \dots, F(x_k))
            \qquad (n \to \infty).
        \]
        \begin{solution}
            Since \(Y_n\coloneq (Y_n^{(1)}, \dots, Y_n^{(k)}) \coloneq (\ind{X_n
            \le x_1}, \dots, \ind{X_n \le x_k})\) are independent and
            identically distributed with expectation \((F(x_1), \dots,
            F(x_k))\), we have by the strong law of large numbers that
            \[
                (F_n(x_1), \dots, F_n(x_k)) = \frac1n \sum_{i=1}^n Y_i
                \overset{\as}\to (F(x_1), \dots, F(x_k)).
                \qedhere
            \]
        \end{solution}
    \end{enumerate}
\end{exercise}

\begin{exercise}[Strong law of large numbers in \(L^4\) \Coffeecup\Coffeecup]
    \label{ex: strong law of large numbers in L4}
    Let \((X_n)_{n\in \nat}\) be a sequence of independent and identically
    distributed random variables in \(\cL^4(\real)\). Bound the centered fourth moment
    \(\|\mean{X}_n - \E{X_1}\|_{L^4}^4\) to prove the strong law of large numbers.

    \noindent\textbf{Hint:} Assume \(\E{X_1} = 0\) without loss of generality.

    \begin{solution}
        Assume that the random variables are centered without loss of generality.
        For indices \(i,j,k,l\), where \(i\) is not equal to any of the others
        we have by independence
        \[
            \E{X_i X_j X_k X_l} = \E{X_i}\E{X_jX_k X_l} = 0.
        \]
        So clearly, if there is an index that does not appear twice the expectation is zero.
        This implies
        \[\begin{aligned}
            \norm{\mean{X}_n}_{L^4}^4
            &= \E[\Big]{\abs[\Big]{\frac1n \sum_{i=1}^n X_i}^4}
            \\
            &= \frac1{n^4} \sum_{i,j,k,l=1}^n \E{X_i X_j X_k X_l}.
            & (\abs{x}^4=x^4)
            \\
            &= \frac1{n^4} \Bigl(
                \sum_{i=1}^n \E{X_i^4}
                + \sum_{\substack{i,j=1\\i\neq j}}^n \E{X_i^2}\E{X_j^2}
            \Bigr)
            & (\text{no single index})
            \\
            &=
            \frac1{n^4}\Bigl(n\E{X_1^4} + \frac{n(n-1)}{2} (\E{X_1^2})^2\Bigr)
            &(\text{identical distributed})
            \\
            &\le \frac{C}{n^2}.
        \end{aligned}
        \]
        With \(C = \E{X_1^4} + \frac12 (\E{X_1^2})^2 < \infty\). Consequently, we have
        by Chebyshev's inequality for \(p=4\) (Corollary \ref{cor: markov inequality Lp})
        \[
            \prob{\norm{\mean{X}_n} > \epsilon} \le \frac{\norm{\mean{X}_n}_{L^4}^4}{\epsilon^4} \le \frac{C}{\epsilon^4 n^2}.
        \]
        Since this is summable in \(n\) we have by Corollary \ref{cor:
        characterization of almost sure convergence} that \(\mean{X}_n
        \overset{\as}\to 0\).
    \end{solution}
\end{exercise}

\begin{exercise}[Variance estimation \Coffeecup]
    \label{ex: variance estimation}
    Let \((X_n)_{n\in \nat}\) be a sequence of \(\iid\) random variables in \(\real\)
    with expectation \(\mu\) and variance \(\sigma^2\). Prove that
    \begin{enumerate}[label=(\alph*),noitemsep]
        \item 
        \(
            \hat\sigma^2_n \coloneq \frac1n\sum_{i=1}^n (X_i-\mu)^2
            \overset{\as}\to \sigma^2,
        \)
        \begin{solution}
            Since \(\E{(X_i-\mu)^2} = \var{X_i} = \sigma^2\), this is
            simply the strong law of large numbers for the independent and identically
            distributed random variables \((X_i-\mu)^2\).
        \end{solution}

        \item Without knowing the expectation, a reasonable variance estimator
        uses the sample mean \(\mean{X}_n\) instead of \(\mu\), that is
        \[
            \hhat\sigma_n^2 \coloneq \frac1n\sum_{i=1}^n (X_i-\mean{X}_n)^2
        \]
        Prove that it converges almost surely to \(\sigma^2\).
        
        \textbf{Hint:} First prove \(\hhat\sigma_n^2 = \hat\sigma_n^2 - (\mean{X}_n - \mu)^2\).
        \begin{solution}
            We have that
            \begin{align}
                \nonumber
                \hhat\sigma_n^2
                &= \frac1n\sum_{i=1}^n (X_i-\mean{X}_n)^2
                \\
                \nonumber
                &= \frac1n\sum_{i=1}^n ((X_i-\mu) - (\mean{X}_n - \mu))^2
                \\
                \nonumber
                &= \frac1n\sum_{i=1}^n (X_i-\mu)^2 - 2(X_i-\mu)(\mean{X}_n - \mu) + (\mean{X}_n - \mu)^2
                \\
                \nonumber
                &= \frac1n\sum_{i=1}^n (X_i-\mu)^2 - 2(\mean{X}_n - \mu)^2 + (\mean{X}_n - \mu)^2
                \\
                \label{eq: variance estimator decomposition}
                &= \hat\sigma_n^2 - (\mean{X}_n - \mu)^2.
            \end{align}
            By \(\mean{X}_n \overset{\as}\to \mu\) and continuous mapping
            we have \((\mean{X}_n - \mu)^2 \overset{\as}\to 0\). Consequently,
            by \(\hat\sigma_n^2 \overset{\as}\to \sigma^2\), joint convergence
            (Theorem \ref{thm: joint convergence}) and continuous mapping we have
            \(\hhat\sigma_n^2 \overset{\as}\to \sigma^2\).
        \end{solution}

        \item Prove that \(\hhat\sigma_n^2\) is biased. Specifically, prove
        \[
            \E{\hhat\sigma_n^2} = \tfrac{n-1}{n}\sigma^2.
        \]
        Prove the consistency of the unbiased estimator \(\tilde\sigma_n^2 \coloneq \tfrac{n}{n-1}\hhat\sigma_n^2\),
        that is \(\tilde\sigma_n^2 \overset{\as}\to \sigma^2\).
        \begin{solution}
            Bienaymé (Lemma \ref{lem: Bienaymé}) implies \(\var{\mean{X}_n} =
            \frac1n \var{X_1}\) and therefore using \eqref{eq: variance estimator decomposition}
            and the linearity of expectation we have
            \begin{align*}
                \E{\hhat\sigma_n^2}
                &= \E{\hat\sigma_n^2} - \E{(\mean{X}_n - \mu)^2}
                \\
                &= \sigma^2 - \var{\mean{X}_n}
                = \sigma^2 - \tfrac1n\var{X_1} = \tfrac{n-1}{n}\sigma^2.
            \end{align*}
            Consequently, \(\hhat\sigma_n^2\) is a biased estimator of
            \(\sigma^2\) and \(\tilde\sigma_n^2\) is an unbiased estimator of
            \(\sigma^2\) by linearity of the expectation.
            Since \(\frac{n}{n-1} \to 1\) and \(\hhat\sigma_n^2 \overset{\as}\to
            \sigma^2\) we have by joint convergence and continuous mapping that
            \(\tilde\sigma_n^2 \overset{\as}\to \sigma^2\).
        \end{solution}

        \item For every \(\sigma_n^2 \in \{\hat\sigma_n^2, \hhat\sigma_n^2,
        \tilde\sigma_n^2\}\) prove that \(\sqrt{\sigma_n^2}\) is a consistent
        estimator of the standard deviation \(\sigma\), that is 
        \(\sqrt{\sigma_n^2} \overset{\as}\to \sigma\).
        \begin{solution}
            Since \(\sqrt{x}\) is continuous for \(x \ge 0\) and \(\sigma_n^2,
            \sigma^2 \ge 0\) the claim follows from continuous mapping.
        \end{solution}
    \end{enumerate}
\end{exercise}

\begin{exercise}[Random vectors in high dimension \Coffeecup]
    In this exercise we will see that high-dimensional spaces have some counter-intuitive
    properties. Let \((X^n)_{n\in \nat}\) and \((Y^n)_{n\in \nat}\) be two
    sequence of independent identically distributed random variables with mean
    zero and finite variance. Let \(X_\dims \coloneq (X^1, \ldots, X^\dims)\)
    and \(Y_\dims \coloneq (Y^1, \ldots, Y^\dims)\) be the corresponding random vectors in
    \(\real^\dims\).
    Prove that
    \begin{enumerate}[label=(\alph*), noitemsep]
        \item asymptotically and in relative terms both vectors have the same
        length in high dimension, that is
        \[
            \frac{\norm{X_\dims}^2}{\norm{Y_\dims}^2} \overset\as\to 1 \qquad (\dims \to \infty).
        \]
        \item their cosine angle is zero asymptotically, that is
        \[
            \frac{\scp{X_\dims, Y_\dims}}{\norm{X_\dims}\norm{Y_\dims}} \overset\as\to 0 \qquad (\dims \to \infty).
        \]
        In other words: random vectors are asymptotically orthogonal in high dimension.
    \end{enumerate}
\end{exercise}

\begin{exercise}[Products of random variables \Coffeecup]
    This entire chapter is about sums of random variables. In this exercise we want
    to demonstrate that these results can be easily translated to products of
    random variables with the help of the logarithm. For
    \((X_n)_{n\in \nat}\) independent and identically distributed random variables
    with \(X_n > 0\) and \(\E{|\ln(X_1)|} < \infty\) prove that the
    `geometric mean' converges almost surely to the `geometric expectation',
    that is'
    \[
        \Bigl(\prod_{i=1}^n X_i \Bigr)^\frac1n
        \overset{\as}\to \exp\bigl(\E[\big]{\ln(X_1)}\bigr).
    \]

    \begin{solution}
        Observe that the logarithm of the geometric mean converges by the law of
        large numbers to the expectation of the logarithm, that is
        \[
            \ln\Bigl(\Bigl(\prod_{i=1}^n X_i \Bigr)^\frac1n\Bigr)
            = \frac1n \sum_{i=1}^n \ln(X_i)
            \overset{\as}\to \E[\big]{\ln(X_1)}.
        \]
        Since \(\exp(x)\) is continuous the claim follows from continuous mapping.
    \end{solution}

\end{exercise}

\paragraph*{Finance application}

In the following two exercises we want to model a financial market with \(n\)
assets. Each asset \(i\) has returns
\begin{equation}
    \label{eq: return definition}    
    R_t^i = \tfrac{P_{t+1}^i}{P_t^i},
\end{equation}
where \(P_t^i\) is the price of the asset \(i\) at time \(t\) and
\(P_{t+1}^i\) is the price of the asset at time \(t+1\).

\begin{exercise}[Long term investment returns \Coffeecup]
    In this exercise we restrict ourselves to a single asset.
    That is \(n=1\) and we write \(R_t = R_t^1\) for the return of the asset at
    time \(t\), cf.~\eqref{eq: return definition}.
    Assume that the returns \(R_0, R_1, \dots\) are independent and identically
    distributed with \(\E{|\ln(R_1)|} < \infty\). Compare the
    geometric mean with the arithmetic mean, that is
    \[
        \Bigl(\prod_{t=0}^{n-1} R_t\Bigr)^\frac1n
        \quad \text{vs.} \quad
        \frac1n\sum_{t=0}^{n-1} R_t.
    \]
    Which one captures the long term growth of the investment better?
    \begin{solution}
        The geometric mean of the returns \(\mean{R}^g_n \coloneq \Bigl(\prod_{t=0}^{n-1} R_t\Bigr)^\frac1n\)
        repeated for \(n\) periods gives the same return as the actual return over \(n\) periods:
        \[
            (\mean{R}^g_n)^n = \prod_{t=0}^{n-1} R_t = \frac{P_n}{P_0}.
        \]
        Consequently, \(r^g \coloneq \exp(\E{\ln(R_1)})\) is the long term
        growth rate that stems from the law of large numbers
        \[
            \mean{R}^g_n  = \exp\Bigl(\frac1n \sum_{t=0}^{n-1} \ln(R_t)\Bigr) \overset{\as}\to r^g.
        \]
        In contrast, \(\E{R_1}\) may be arbitrarily far from \(r^g\) and
        \(\E{R_1}^n\) is thereby not a good approximation of the return over \(n\) periods.
        The arithmetic mean approximates \(\E{R_1}\) and is thereby not a good
        approximation of the long term growth of the investment.
    \end{solution}
\end{exercise}

\begin{exercise}[Diversification and systemic risk \Coffeecup\Coffeecup]
    In this exercise we will only model a single time step from \(t=0\) to
    \(t=1\) and write \(R^i = R_0^i\), cf.~\eqref{eq: return definition}. Moreover we assume for \(\alpha^i,
    \beta^i \in \real\)
    \[
        R^i = 1
        + \underbrace{\alpha^i}_{\substack{\text{expected}\\ \text{return}}}
        + \underbrace{\beta^i M}_{\substack{\text{systemic}\\\text{risk}}}
        + \underbrace{\epsilon^i}_{\substack{\text{specific}\\\text{risk}}},
    \]
    where \(M\) is a random variable modelling the ``market risk'' or ``systemic
    risk'' and \(\epsilon^i\) is a random variable modelling the ``specific
    risk'' of the asset \(i\). We assume that the random variables \(M\) and
    \(\epsilon^1, \dots, \epsilon^n\) are uncorrelated with zero mean. Without
    loss of generality we assume that the variance of \(M\) is one -- simply move
    the variance into \(\beta^i\). We also assume that the variance of
    \(\epsilon^i\) is bounded by \(\sigma_\epsilon^2\). For simplicity we
    further assume \(\alpha^i = \alpha\) and \(\beta^i = \beta\) for all \(i\).

    We consider the portfolio \(P\) that invests equally in all assets such that
    the portfolio return is given by
    \[
        R^P = \frac1n \sum_{i=1}^n R^i.
    \]
    Prove that
    \begin{enumerate}[label=(\alph*),noitemsep]
        \item The expected return of the portfolio is given by \(\E{R^P} = 1 + \alpha\).
        \begin{solution}
            By linearity of the expectation and the assumptions on the random variables we have 
            \[
                \E{R^P}
                = \frac1n \sum_{i=1}^n \E{R^i}
                = \frac1n \sum_{i=1}^n (1 + \alpha + \beta \E{M} + \E{\epsilon^i})
                = 1 + \alpha.
                \qedhere
            \]
        \end{solution}
        
        \item The variance of the portfolio approaches \(\beta^2\) from above as the
        diversification grows, that is \(n\to \infty\). 
        \textbf{Hint:} Lemma \ref{lem: Bienaymé}.
        \begin{solution}
            To use Bienaymé (Lemma \ref{lem: Bienaymé}) we calculate the
            variance and covariance. The covariance is given by
            \begin{align*}
                \cov{R^i}{R^j}
                &= \E{(R^i - \E{R^i})(R^j - \E{R^j})}
                \\
                &= \E{(\beta M + \epsilon^i)(\beta M + \epsilon^j)}
                \\
                &= \beta^2 \underbrace{\E{M^2}}_{=1} + \beta \underbrace{\E{M\epsilon^i}}_{=0} +
                \beta \underbrace{\E{M\epsilon^j}}_{=0} +
                \E{\epsilon^i\epsilon^j}
                = \beta^2 + \E{\epsilon^i\epsilon^j}.
            \end{align*}
            Consequently, if \(i \neq j\) we have \(\cov{R^i}{R^j} = \beta^2\) and if \(i=j\) we have
            \[
                \var{R^i} = \beta^2 + \var{\epsilon^i} \le \beta^2 + \sigma_\epsilon^2
            \]
            Using Bienaymé we have consequently have
            \begin{align*}
                \abs{\var{R^P} - \beta^2}
                &= \abs[\Big]{\frac1{n^2}\sum_{i=1}^n \var{R^i} + \frac1{n^2}\sum_{i\neq j} \cov{R^i}{R^j} - \beta^2}
                \\
                &\le \tfrac1{n}(\beta^2 + \sigma_\epsilon^2) + \abs{(1-\tfrac1{n})\beta^2 - \beta^2}
                \\
                &= \frac{\beta^2 + \sigma_\epsilon^2 + \beta^2}{n}
                \to 0.
                \qedhere
            \end{align*}
        \end{solution}

    \end{enumerate}
\end{exercise}

\begin{remark}
    Diversification over companies is additive, while the long term view
    is multiplicative. This makes it somewhat delicate to combine both views
    into one. Nevertheless both views reduce the variance and thereby the risk
    of the investment by the law of large numbers.
\end{remark}


\begin{exercise}[Glivenko-Cantelli theorem \Coffeecup\Coffeecup\Coffeecup]
    This exercise is an extension of Exercise \ref{ex: mean of indicators} where
    we want to prove the Glivenko-Cantelli theorem, which says that the
    empirical distribution function \(F_n\) converges uniformly to the
    distribution function and not just at finitely many points.
    \begin{mdframed}[innertopmargin=0pt]
    \begin{theorem}[Glivenko-Cantelli]
        \label{thm: glivenko cantelli}
        Let \((X_n)_{n\in \nat}\) be a sequence of iid random
        variables in \(\real\) with distribution function \(F\) and empirical 
        distribution function \(F_n\). Then
        \[
            \sup_{x\in \real}\abs{F_n(x) - F(x)} \overset{\as}\to 0.
        \]
    \end{theorem}
    \end{mdframed}
    You will prove this theorem in three steps:
    \begin{enumerate}[label=(\alph*),noitemsep]
        \item Define the functions \(G(x) \coloneq \prob{X_1 < x}\) and \(G_n(x) \coloneq \frac1n \sum_{i=1}^n \ind{X_i < x}\)
        and prove for any \(x_1, \dots, x_k \in \real\)
        \[\begin{aligned}
            &(F_n(x_1), \dots, F_n(x_k), G_n(x_1), \dots, G_n(x_k))
            \\
            &\overset{\as}\to (F(x_1), \dots, F(x_k), G(x_1), \dots, G(x_k)).
        \end{aligned}
        \]
        Deduce
        \begin{equation}
            \label{eq: discretization max to zero}            
            R_n \coloneq \max_{i\in \set{1,\dots, k}} \max\set[\big]{\abs{F_n(x_i) - F(x_i)}, \abs{G_n(x_i) - G(x_i)}} \overset{\as}\to 0.
        \end{equation}
        \begin{solution}
            The first claim follows from the law of large numbers as in Exercise \ref{ex: mean of indicators}.
            The claim \eqref{eq: discretization max to zero} then follows from the norm equivalence
            in \(\real^{2k}\), that is there exists a constant \(C\) such that
            \[
                \norm{x}_\infty \le C \norm{x}_2
            \]
            As we have convergence in \(\norm{\cdot}_2\) we thereby also have
            convergence in \(\norm{\cdot}_\infty\) and consequently \(R_n
            \overset{\as}\to 0\).
        \end{solution}

        \item\label{eq: glivenko cantelli limsup bound} Recall the properties of the quantile function \(F^\leftarrow\) (Exercise \ref{ex: stochastic simulation})
        and choose \(x_1, \dots, x_k\) as \(x_i \coloneq F^\leftarrow\bigl(\frac{i}k\bigr)\).
        Prove
        \begin{equation}
            \label{eq: discretization properties}    
            G(x_i) \le \tfrac{i}k \le F(x_i).
        \end{equation}
        Then use this and \eqref{eq: discretization max to zero} to prove that
        \[
            \limsup_{n\to \infty} \sup_{x\in \real}\abs{F_n(x) - F(x)} \le \tfrac1k + \limsup_{n\to \infty}R_n \le \tfrac1k.
        \]
        \begin{solution}
            The upper bound in \eqref{eq: discretization properties} follows from Exercise \ref{ex: stochastic simulation}
            \ref{it: F(F-(p)) > p} and the lower bound follows from the fact that
            \(G(x) \le F(x) \le \frac{i}k\) for all \(x < x_i\) due to
            \[
                x_i = F^\leftarrow\bigl(\tfrac{i}k\bigr)
                = \inf\set[\big]{x\in \real: F(x) \ge \tfrac{i}k}.
            \]
            and left continuity of \(G(x)\) by
            \begin{align*}
                G(x) = \prob{X_1 < x}
                &= \prob[\Big]{\bigcup_{n\in \nat} \set{X_1 < y- \tfrac1n}}
                \\
                &= \lim_{n\to \infty} \prob{X_1 < y} = \lim_{n\to\infty} G(y - \tfrac1n).
            \end{align*}
            On to the second claim: For any \(x\in (x_i, x_{i+1})\) we have
            \begin{align*}
                F_n(x)
                \overset{x< x_{i+1}}&\le G_n(x_{i+1})
                \\
                \overset{\eqref{eq: discretization max to zero}}&\le G(x_{i+1}) + R_n
                \\
                \overset{\eqref{eq: discretization properties}}&\le F(x_i) + \frac1k + R_n
                \overset{x_i < x}\le F(x) + \frac1k + R_n
            \end{align*}
            Similarly, we have for any \(x\in (x_i, x_{i+1})\) that
            \begin{align*}
                F_n(x)
                \overset{x>x_i}&\ge F_n(x_i)
                \\
                \overset{\eqref{eq: discretization max to zero}}&\ge F(x_i) - R_n
                \\
                \overset{\eqref{eq: discretization properties}}&\ge G(x_{i+1}) - \tfrac1k - R_n
                \\
                \overset{x_{i+1} > x}&\ge F(x) - \tfrac1k - R_n.
            \end{align*}
            Put together, we have for any \(x\in (x_i, x_{i+1})\) and any \(i \in \set{1,\dots, k-1}\)
            \begin{equation}
                \label{eq: glivenko cantelli bound}    
                \abs{F_n(x) - F(x)} \le \frac1k + R_n
            \end{equation}
            Since we clearly also have this for for \(x = x_i\) by definition of
            \(R_n\). What is therefore left are only the cases \(x< x_1\) and \(x>x_k\). For
            \(x< x_1\) we have
            \[
                F_n(x) \overset{x < x_1}\le G_n(x_1)
                \overset{\eqref{eq: discretization max to zero}}\le G(x_1) + R_n
                \overset{\eqref{eq: discretization properties}}\le \tfrac1k + R_n
                \overset{F(x) \ge 0}\le F(x) + \tfrac1k + R_n,
            \]
            and
            \[
                F_n(x) \overset{x < x_1}\ge 0 
                \overset{\eqref{eq: discretization properties}}\ge G(x_1) - \tfrac1k
                \overset{x< x_1}\ge F(x) - \tfrac1k.
            \]
            Together these result in \eqref{eq: glivenko cantelli bound} for \(x < x_1\). For \(x > x_k\) we have
            \[
                F_n(x) \overset{x > x_k}\le 1
                \overset{\eqref{eq: discretization properties}}\le F(x_k) \le F(x)
            \]
            and
            \[
                F_n(x) \overset{x > x_k}\ge F_n(x_k)
                \overset{\eqref{eq: discretization max to zero}}\ge F(x_k) - R_n
                \overset{\eqref{eq: discretization properties}}\ge 1 - R_n
                \overset{F(x) \le 1}\ge F(x) - R_n.
            \]
            Together they result in \eqref{eq: glivenko cantelli bound} for \(x > x_k\). In summary, we have \eqref{eq: glivenko cantelli bound} 
            for all \(x\in \real\) and therefore almost surely
            \[
                \limsup_{n\to \infty} \sup_{x\in \real}\abs{F_n(x) - F(x)} \le \tfrac1k + \limsup_{n\to \infty}R_n \overset{\eqref{eq: discretization max to zero}}= \tfrac1k.
                \qedhere
            \]
        \end{solution}

        \item Deduce the ``Glivenko-Cantelli theorem'' with a union bound over \(k\).
        \begin{solution}
            We have
            \begin{align*}
                &\prob[\Big]{\limsup_{n\to \infty} \sup_{x\in \real}\abs{F_n(x) - F(x)} > 0}
                \\
                &= \prob[\Big]{\bigcup_{k\in \nat}\limsup_{n\to \infty} \sup_{x\in \real}\abs{F_n(x) - F(x)} > \tfrac1k}
                \\
                &= \lim_{k\to\infty}\prob[\Big]{\limsup_{n\to \infty} \sup_{x\in \real}\abs{F_n(x) - F(x)} > \tfrac1k}
                \\
                \overset{\ref{eq: glivenko cantelli limsup bound}}&= \lim_{k\to \infty} 0 =0.
            \end{align*}
            and thereby
            \begin{align*}
                &\prob[\Big]{\lim_{n\to \infty} \sup_{x\in \real}\abs{F_n(x) - F(x)} = 0}
                \\
                &= 1- \prob[\Big]{\limsup_{n\to \infty} \sup_{x\in \real}\abs{F_n(x) - F(x)} > 0}
                \\
                &= 1.
                \qedhere
            \end{align*}
        \end{solution}
    \end{enumerate}
\end{exercise}
